\section{PCDVQ}

\begin{frame}{Polar Coordinate Decoupling Vector Quantization (PCDVQ)}
    \begin{columns}
        \column{0.55\textwidth}
            \centering
            \includegraphics[width=\linewidth,height=0.34\textheight,keepaspectratio]{pipeline.png}
            \captionof{figure}{The pipeline of PCDVQ}\label{fig:pipeline}
            \vspace{0.5em}

            \begin{minipage}{\linewidth}
                \centering
                \includegraphics[width=0.49\linewidth,height=0.4\textheight,keepaspectratio]{sensitivity_gap.png}\hfill
                \includegraphics[width=0.49\linewidth,height=0.4\textheight,keepaspectratio]{error_diff.png}\\
                {\captionof{figure}{Compare direction and magnitude}\label{fig:compare_dir_magn}}
            \end{minipage}

        \column{0.45\textwidth}
            This method consist of the next stages (Fig.~\ref{fig:pipeline}):
            \begin{itemize}
                \item Standard Gaussian Regularization (SGR). They used Random Hadamard Transformation (RHT) for SGR
                \item Polar Coordinate Decoupling (PCD)
                \item Distribution Aligned Codebook Construction (DACC)
            \end{itemize}

            Motivation: direction is identified to be more essential than magnitude (Fig.~\ref{fig:compare_dir_magn})
    \end{columns}
\end{frame}
